\documentclass[aspectratio=169]{beamer}

% EBU6505 module-introduction deck for Semester 1, 2026--2027.
% Academic profile details verified against https://riasatislam.com/experience/.
\usepackage[utf8]{inputenc}
\usepackage{amsmath,amssymb}
\usepackage{booktabs}
\usepackage{tabularx}
\usepackage{tcolorbox}
\usepackage{hyperref}

\usetheme{Madrid}
\usecolortheme{default}
\setbeamertemplate{navigation symbols}{}
\setbeamertemplate{footline}{%
  \leavevmode%
  \hbox{%
    \begin{beamercolorbox}[wd=.70\paperwidth,ht=2.5ex,dp=1ex,leftskip=.4em]{author in head/foot}
      \scriptsize EBU6505 Reasoning and Agents
    \end{beamercolorbox}%
    \begin{beamercolorbox}[wd=.30\paperwidth,ht=2.5ex,dp=1ex,rightskip=.4em]{date in head/foot}
      \scriptsize\hfill Semester 1, 2026--2027
    \end{beamercolorbox}%
  }%
}

\title[EBU6505 Module Introduction]{EBU6505 -- Reasoning and Agents\\Module Introduction}
\author{Dr Riasat Islam}
\institute{Queen Mary University of London}
\date{Semester 1, 2026--2027}

\begin{document}

\begin{frame}
  \titlepage
\end{frame}

\begin{frame}{Teaching Team}
\begin{columns}[T]
  \begin{column}{0.48\textwidth}
    \begin{block}{Dr Riasat Islam}
    Lecturer in Computer Science\\
    Programme Director, Intelligent Science and Technology

    \vspace{0.15cm}
    \textbf{Teaching:} Block 1 (Teaching Week 3) and Block 2 (Teaching Week 4)
    \end{block}
  \end{column}
  \begin{column}{0.48\textwidth}
    \begin{block}{Dr Chao Shu}
    Lecturer in Telecommunication Engineering\\
    Programme Director, Telecommunications Engineering with Management

    \vspace{0.15cm}
    \textbf{Teaching:} Block 3 (Teaching Week 9) and Block 4 (Teaching Week 10)
    \end{block}
  \end{column}
\end{columns}

\vspace{0.4cm}
\begin{center}
  \textbf{We will connect formal ideas to practical systems, current AI tools, and responsible decision-making.}
\end{center}
\end{frame}

\begin{frame}{Meet Your Block 1 Lecturer}
\begin{block}{Dr Riasat Islam}
I joined Queen Mary University of London as a Lecturer in August 2024.
\end{block}

\begin{itemize}
  \item \textbf{BSc Engg in Electrical and Electronic Engineering}, Islamic University of Technology, Bangladesh (2009--2012)
  \item \textbf{MSc in ICT Innovation}, University College London and KTH Royal Institute of Technology (2013--2015)
  \item \textbf{PhD in Computing}, The Open University (2017--2023)
  \item Interests spanning artificial intelligence, human--computer interaction, healthcare, well-being, and education
\end{itemize}

\vspace{0.2cm}
\begin{center}
\small\url{https://riasatislam.com}
\end{center}
\end{frame}

\begin{frame}{Why This Module Now?}
\begin{block}{AI is rapidly changing}
We hope this module helps you understand the underlying concepts and equips you with the skills to navigate this exciting era of AI.
\end{block}

\begin{itemize}
  \item Move beyond using AI tools to understanding how reasoning systems work.
  \item Learn to question assumptions, inspect evidence, and evaluate outputs.
  \item Connect algorithms to real systems, people, organisations, and consequences.
\end{itemize}

\begin{center}
\Large\textbf{Let’s have fun!}
\end{center}
\end{frame}

\begin{frame}{What You Will Learn}
\begin{itemize}
  \item Model problems using states, actions, knowledge, uncertainty, and goals.
  \item Apply reasoning methods to solve problems and explain decisions.
  \item Build and critique knowledge-grounded and probabilistic agents.
  \item Use modern AI tools thoughtfully, with verification and reflection.
  \item Design systems that are useful, transparent, robust, and responsible.
\end{itemize}

\vspace{0.25cm}
\begin{center}
\textbf{The aim is not only to get an answer, but to understand why an answer should be trusted.}
\end{center}
\end{frame}

\begin{frame}{Block Overview}
\scriptsize
\renewcommand{\arraystretch}{1.08}
\begin{tabularx}{\textwidth}{>{\bfseries}l X X}
\toprule
Block & Teaching week & Main themes \\
\midrule
1 & 3 & Reinforcement learning; knowledge representation; knowledge graphs; modern coding agents \\
\addlinespace[8pt]
2 & 4 & Bayesian networks; exact inference; probabilistic reasoning over time; planning under uncertainty; DDNs and POMDPs \\
\addlinespace[8pt]
3 & 9 & Further agent architectures and reasoning topics led by Dr Chao Shu \\
\addlinespace[8pt]
4 & 10 & Further decision-making under uncertainty topics led by Dr Chao Shu \\
\bottomrule
\end{tabularx}

\vspace{0.15cm}
\begin{center}\scriptsize
Blocks 1 and 2 establish the reasoning foundations; Blocks 3 and 4 extend them to wider agent and decision-making systems.
\end{center}
\end{frame}

\begin{frame}{Teaching Schedule}
\begin{center}
\renewcommand{\arraystretch}{1.2}
\begin{tabular}{>{\bfseries}l l l}
\toprule
Teaching week & Block & Lecturer \\
\midrule
3 & Block 1 & Dr Riasat Islam \\
\addlinespace[8pt]
4 & Block 2 & Dr Riasat Islam \\
\addlinespace[8pt]
7--8 & Coursework Mini Project 1 & 20\% \\
\addlinespace[8pt]
9 & Block 3 & Dr Chao Shu \\
\addlinespace[8pt]
10 & Block 4 & Dr Chao Shu \\
\addlinespace[8pt]
11--12 & Coursework Mini Project 2 & 20\% \\
\addlinespace[8pt]
After week 16 & Exam (date to be confirmed) & 60\% \\
\bottomrule
\end{tabular}
\end{center}
\end{frame}

\begin{frame}{How We Will Learn}
\begin{block}{Problem-solving based learning}
We will work through concrete problems, compare possible approaches, explain our reasoning, and learn from mistakes.
\end{block}

\begin{block}{AI-assisted learning: support, not substitution}
AI can support your learning, but it does not replace your own reasoning. Start by thinking through the problem yourself, then use AI as a tutor, critic, or debugging partner.
\end{block}

\begin{columns}[T]
  \begin{column}{0.48\textwidth}
    \scriptsize
    \textbf{Useful learning activities}
    \begin{itemize}\scriptsize\setlength{\itemsep}{0pt}
      \item Ask for hints, explanations, alternatives, or counterexamples.
      \item Ask the tool to challenge assumptions, test edge cases, and debug code.
    \end{itemize}
  \end{column}
  \begin{column}{0.48\textwidth}
    \scriptsize
    \textbf{Your responsibility}
    \begin{itemize}\scriptsize\setlength{\itemsep}{0pt}
      \item Verify outputs against lectures and reliable references.
      \item Explain every step in your own words; AI can be fluent, incomplete, or wrong.
      \item Follow each assessment brief: never outsource assessed reasoning or present generated work as your own.
    \end{itemize}
  \end{column}
\end{columns}

\vspace{-0.05cm}
\begin{center}\scriptsize
Students who attend regularly, attempt the problems, and participate actively tend to perform very well.
\end{center}
\end{frame}

\begin{frame}{Assessment}
\begin{columns}[T]
  \begin{column}{0.47\textwidth}
    \begin{block}{Coursework -- 40\%}
    \begin{itemize}
      \item Mini Project 1: 20\%\\
      \small Teaching Weeks 7--8
      \item Mini Project 2: 20\%\\
      \small Teaching Weeks 11--12
    \end{itemize}
    \end{block}
  \end{column}
  \begin{column}{0.47\textwidth}
    \begin{block}{Exam -- 60\%}
    Held after Teaching Week 16.\\[0.2cm]
    Date and further arrangements will be confirmed later.
    \end{block}
  \end{column}
\end{columns}

\vspace{0.35cm}
\begin{center}
Coursework rewards sustained problem-solving across the module; the exam assesses your understanding and application of the core ideas.
\end{center}
\end{frame}

\begin{frame}{Module Representatives}
\begin{center}
\scriptsize EBU6505 \hfill riasat.islam@qmul.ac.uk \hfill Semester 1, 2026--2027
\end{center}
\vspace{0.25cm}
\begin{itemize}
  \item Collect feedback from classmates about the lectures.
  \item Assist lecturers in disseminating important information.
  \item Attend one Student--Staff Liaison Committee (SSLC) meeting during the semester.
  \item {\color{red}\textbf{Volunteers (one or two) will be welcomed or chosen during the first teaching week of the module.}}
\end{itemize}

\vspace{0.45cm}
\begin{center}
\textbf{Module representatives help us listen, communicate, and improve the module together.}
\end{center}
\end{frame}

\end{document}
